\uuid{dOIV}
\titre{Intervalle}
\auteur{bodin ; hanani ; mzari}
\organisation{exo7}
\licence{CC BY-NC-SA 4.0}
\datecreate{2020-01-20}
\niveau{L1}
\module{Analyse}
\chapitre{Propriétés de R}
\sousChapitre{Intervalle, densité}
\difficulte{3}
\theme{Intervalle, densité, intervalle}
\exo7qcmid{lille-1.reels.026}

\contenu{
\begin{question}
\qclassification{120.04}
\qtags{motcle=intervalle}
Soient $I$ et $J$ deux intervalles de $\Rr$. Quelles sont les assertions vraies ?
\begin{answers}
    \bad{$I \cup J$ est un intervalle.}

    \good{$I \cap J$ est un intervalle (éventuellement réduit à un point ou vide).}

    \good{Si $I \cap J \neq \varnothing$ alors $I \cup J$ est un intervalle.}

    \good{Si $I \subset J$ alors $I \cup J$ est un intervalle.}
\end{answers}
\begin{explanations}
Tracer les intervalles sur la droite réelle pour mieux comprendre.
Une union d'intervalles n'est en général pas un intervalle !
\end{explanations}
\end{question}
}
